
\textbf{Актуальность темы} определяется тем, что для современных edge/IoT-сред
одновременно характерны высокая динамичность, распределённый характер
управления и ограниченность локальных вычислительных ресурсов. В таких
условиях методы, основанные только на детерминированной оптимизации, требуют
слишком сильных предположений о полноте и стабильности информации о системе,
тогда как purely learning-based подходы, в том числе RL и MARL, хотя и обладают
адаптивностью, обычно не содержат встроенного механизма экономически
интерпретируемого согласования интересов участников.

Отдельную трудность представляет то, что в распределённой edge-среде решения
принимаются многократно и последовательно: текущее назначение задач влияет на
будущие очереди, загрузку узлов, сетевые задержки и последующие возможности
распределения ресурсов. Следовательно, задача должна рассматриваться не как
одноразовая оптимизация, а как последовательный процесс управления в
многоагентной частично наблюдаемой среде.

При этом в литературе сохраняется разрыв между двумя линиями исследований. С одной стороны, \textbf{аукционные и механизм-дизайн подходы} позволяют формально описывать стимулы, платежи, правдивость заявок и максимизацию общественного благосостояния, однако в классическом виде хуже приспособлены к динамической нестационарной среде. С другой стороны, \textbf{кооперативное многоагентное обучение с подкреплением} позволяет строить адаптивные политики в сложной среде, но само по себе не обеспечивает явного механизма стимулирования и интерпретируемого учёта внешних эффектов перераспределения ресурсов.

Тем самым актуальной является задача синтеза такого подхода, который сочетал бы экономически осмысленный механизм распределения, строгую многоагентную формализацию и адаптивную обучаемую политику, пригодную для децентрализованного исполнения в edge-среде.

\textbf{Целью} настоящей диссертационной работы является разработка моделей и гибридного метода адаптивного распределения вычислительных ресурсов в распределённых многоагентных edge-инфраструктурах в условиях частичной наблюдаемости и динамической нагрузки на основе интеграции аукционного механизма VCG и кооперативного многоагентного обучения с подкреплением.

Для достижения поставленной цели необходимо решить следующие \textbf{задачи}.
\begin{enumerate}
	\item Провести анализ существующих методов распределения ресурсов в edge- и IoT-системах, выявить ограничения детерминированных, аукционных и RL/MARL-подходов и обосновать требования к разрабатываемому методу;
	\item Разработать математическую модель распределения ресурсов в распределённой edge-системе с учётом гетерогенности ресурсов, ограничений по задержке, характеристик задач и экономических параметров взаимодействия агентов;
	\item Формализовать задачу адаптивного управления ресурсами в виде расширенной модели Dec-POMDP и обосновать применение кооперативного MARL-подхода в парадигме централизованного обучения и децентрализованного исполнения;
	\item Разработать гибридный метод и алгоритм MA--VCG--QMIX, интегрирующий аукционный механизм распределения ресурсов и обучаемую многоагентную политику принятия решений;
	\item Провести экспериментальную оценку предложенного метода в ряде сценариев и определить его преимущества, ограничения и области применимости.
	\item Определить условия практического применения, ограничения, архитектурное место и переносимость предложенного метода для прикладных edge-сценариев, включая роботизированные складские системы, промышленный AIoT и интеллектуальные транспортные подсистемы.
\end{enumerate}

\textbf{Научная новизна} диссертации состоит в следующем:
\begin{enumerate}
	\item Предложена расширенная модель распределения вычислительных ресурсов в edge-системе, объединяющая ресурсные ограничения, сетевые задержки, характеристики задач и экономические параметры взаимодействия агентов;
	\item Предложена формализация задачи адаптивного распределения ресурсов в виде расширенной модели Dec-POMDP, в которой экономические сигналы аукционного механизма включены в контур принятия решений агентов;
	\item Разработан гибридный метод MA--VCG--QMIX, сочетающий VCG-подобный аукционный слой и кооперативное многоагентное обучение с подкреплением в парадигме CTDE;
	\item Предложен способ включения аукционных сигналов в пространство состояний, наблюдений, допустимых действий и вознаграждения обучаемой многоагентной системы;
	\item Разработана программа экспериментальной валидации предложенного метода, обеспечивающая проверку его эффективности, адаптивности, устойчивости и fairness-свойств на воспроизводимом наборе сценариев.
	\item Экспериментально установлены условия применимости предложенного гибридного подхода в сценариях стабильной, пиковой и нарушенной нагрузки, а также выявлен компромисс между эффективностью распределения, справедливостью и адаптивностью.
\end{enumerate}

\textbf{Теоретическая значимость} работы состоит в развитии формализации задач распределения ресурсов в многоагентных edge-системах с частичной наблюдаемостью, а также в исследовании возможностей интеграции механизмов стимулирования и кооперативного MARL в едином контуре управления.

\textbf{Практическая значимость} работы состоит в возможности применения разработанного метода для построения систем распределения вычислительных задач в edge-инфраструктурах с динамической нагрузкой и ограниченной наблюдаемостью, в частности в роботизированных складских системах, промышленном AIoT, распределённой робототехнике и интеллектуальных транспортных подсистемах.

\textbf{Обоснованность и достоверность} полученных результатов обеспечивается:
\begin{enumerate}
	\item использованием признанных теоретических оснований: теории игр, аукционной теории, теории марковских процессов принятия решений, Dec-POMDP и методов MARL;
	\item согласованностью математической модели распределения ресурсов, формальной многоагентной постановки и гибридного алгоритмического решения;
	\item воспроизводимым протоколом экспериментальной валидации, включающим сценарное тестирование, сравнение с baseline-методами, многократные независимые запуски и статистическую обработку результатов;
	\item возможностью независимой проверки алгоритма на реализации в репозитории с использованием фиксированных конфигураций среды и параметров обучения.
\end{enumerate}

\textbf{Апробация работы.} Результаты, изложенные в данной диссертации, докладывались и обсуждались на конференциях:
\begin{enumerate}
	\item 11-я Международная конференция <<Распределенные вычисления и грид-технологии в науке и образовании>> (GRID'2025), Российская Федерация, г. Дубна, 07.07.2025–11.07.2025;
	\item Ежегодная всероссийская научная конференция по проблемам информатики в секции <<Методы машинного обучения и архитектуры интеллектуальных систем>>, Российская Федерация, г. Санкт-Петербург, 28.04.2026.
\end{enumerate}

\textbf{Публикации.} Основные результаты по теме диссертации изложены в 4 печатных изданиях, из которых ... — в периодических научных журналах, индексируемых Scopus [...], ... — в периодических научных журналах, индексируемых РИНЦ [...]. Получены 1 свидетельства о государственной регистрации программ для ЭВМ [...]. TODO 

\textbf{Основные положения, выносимые на защиту:}
\begin{enumerate}
	\item Математическая модель распределения вычислительных ресурсов в распределенных многоагентных системах периферийных вычислений, учитывающая гетерогенность узлов, ограничения по вычислительным ресурсам и задержкам, а также экономические характеристики взаимодействия агентов.
	\item Формализация задачи адаптивного управления распределением ресурсов в виде расширенной модели Dec-POMDP с частичной наблюдаемостью и стохастической динамикой среды.
	\item Гибридный метод MA-VCG-QMIX, объединяющий аукционное распределение ресурсов по механизму VCG и кооперативное многоагентное обучение с подкреплением в парадигме CTDE.
	\item Способ включения экономических сигналов аукционного механизма в контур обучения агентов, обеспечивающий совместный учет эффективности распределения, адаптивности и справедливости.
	\item Программа экспериментальной валидации и бенчмаркинга, предназначенная для воспроизводимой проверки эффективности, устойчивости и границ применимости предложенного метода.
	\item Результаты экспериментального исследования, показывающие, что в сценариях стабильной и пиковой нагрузки предложенный метод улучшает показатели распределения ресурсов по сравнению с базовыми подходами, а в сценариях высокой неоднородности и отказов выявляются границы его применимости.
\end{enumerate}

\textbf{Объем и структура работы.} Диссертация состоит из~введения,
\formbytotal{totalchapter}{глав}{ы}{}{},
заключения и
\formbytotal{totalappendix}{приложен}{ия}{ий}{}.
%% на случай ошибок оставляю исходный кусок на месте, закомментированным
%Полный объём диссертации составляет  \ref*{TotPages}~страницу
%с~\totalfigures{}~рисунками и~\totaltables{}~таблицами. Список литературы
%содержит \total{citenum}~наименований.
%
Полный объём диссертации составляет
\formbytotal{TotPages}{страниц}{у}{ы}{}, включая
\formbytotal{totalcount@figure}{рисун}{ок}{ка}{ков} и
\formbytotal{totalcount@table}{таблиц}{у}{ы}{}.
Список литературы содержит
\formbytotal{citenum}{наименован}{ие}{ия}{ий}.

\textbf{Во введении} сформулированы критерии, показана актуальность и новизна исследования, описана теоретическая и практическая значимость полученных результатов, обозначена цель и задачи исследования для ее достижения, сформулированы положения, выносимые на защиту.

\textbf{В первой главе} проведен анализ существующих подходов к управлению вычислительными ресурсами в распределенных системах периферийных вычислений. Рассмотрены детерминированные, аукционные и многоагентные методы, а также подходы на основе обучения с подкреплением. Выявлены их ограничения и обоснована необходимость разработки гибридного метода.

\textbf{Во второй главе} разработана математическая модель распределения ресурсов в распределенной edge-системе. Сформулированы модели задач и ресурсов, введены функции полезности и стоимости, а также формализована задача максимизации социального благосостояния и описан аукционный механизм распределения ресурсов.

\textbf{В третьей главе} предложена модель адаптивного управления ресурсами на основе многоагентного обучения с подкреплением. Задача формализована в виде Dec-POMDP, описаны состояния, наблюдения, действия, функция вознаграждения и обоснован выбор алгоритма QMIX.

\textbf{В четвертой главе} разработан гибридный метод управления ресурсами, объединяющий аукционный механизм VCG и алгоритм QMIX. Предложена расширенная модель Dec-POMDP, интегрирующая экономические сигналы, и сформулирован алгоритм MA--VCG--QMIX.

\textbf{В пятой глава} приведено сформулирована программа экспериментальной валидации и
бенчмаркинга предложенного метода. Программа включает основные элементы формальной модели: устройства, edge-узлы, сетевую топологию, поток задач, аукционное распределение, VCG-платежи, reward-сигнал и обучение координированной политики.

\textbf{В шестой главе} проведена экспериментальная оценка эффективности разработанного метода. Рассмотрены различные сценарии работы системы, выполнено сравнение с базовыми подходами и проанализированы полученные результаты.

\textbf{В седьмой главе } рассмотрены практические аспекты применения предложенного метода, его место в прикладной архитектуре, условия внедрения, ограничения и переносимость на роботизированные складские системы, промышленный AIoT, интеллектуальные транспортные подсистемы и другие edge-сценарии.

\textbf{В заключении} подведены итоги исследования, обсуждены ограничения, определены направления дальнейшей работы.
